\section{Discussion and limitations}

The risk in \cref{obj:def:curve-risk} measures expected squared supremum error over the budget interval \(I\). Thus, \cref{obj:thm:curve-upper} gives a simultaneous guarantee for one estimated value curve. When \(b_0=0\), that guarantee covers every budget from zero to full population capacity. Together with \cref{obj:thm:capacity-active-lower}, the upper bound identifies the minimax rate under the stated conditions.

When \(d/[n\log(ed)]\) is small, the minimax risk has that order; when the ratio is large, it has constant order. The minimax rate contains a logarithmic factor, and the upper bound uses polynomial estimation of the discrete threshold process. In the paired family, differences among positive cell benefits determine their ranking and hence the half-budget allocation value, even as the full-budget value stays fixed. Under \cref{obj:thm:capacity-active-lower}, this single coordinate attains the curve’s risk order. The factor-one dual contraction carries the process estimate to all budgets, and the whole curve attains the unrestricted full-capacity scalar minimax order.

\subsection{Limitations and future work}

The global supremum criterion covers the endpoint when \(b_0=0\); a budget-local analysis could characterize risk as the budget shrinks with sample size. Adaptive simultaneous bands call for a coverage criterion, while policy regret calls for evaluating an estimated allocation. The theorem's estimator is the exact conditional average in \cref{obj:def:jf-estimator}; approximating that average by Monte Carlo draws would require a separate draw-count and approximation-error guarantee before claiming the same rate. Other directions include covariate alphabets learned from data and allocation across multiple treatment arms.
